\section{場の量子論の基本枠組み}

\noindent\textbf{この章の中心命題}
\[
\boxed{
\mathcal{H}=\mathcal{F}_{\pm}(\mathfrak{h}),
\qquad
[\hat{\phi}(\mathbf{x},t),\hat{\pi}(\mathbf{y},t)]
=i\delta^{(3)}(\mathbf{x}-\mathbf{y})
}
\]
場の量子論では、固定粒子数の一つの空間ではなくフォック空間を状態空間とし、場とその共役運動量を演算子として量子化する。

ここで、これまで別々に準備してきた二つの流れが合流する。ひとつはハイゼンベルク描像で見た「物理量を演算子として記述する」という流れであり、もうひとつは相対論的理論が要請する「粒子数が変化しうる」という流れである。場の量子論では、状態空間は固定粒子数のヒルベルト空間の直和として与えられ、基本変数は時空点に付随した場演算子になる。

\subsection{粒子ではなく場を量子化する}
問題はもはや「よりよい一粒子方程式を探すこと」ではなく、「粒子数が変化しうる理論をどのような状態空間と演算子で定式化するか」にある。相対論的な一粒子波動方程式だけでは、負エネルギー解や反粒子の存在を無理なく扱うことができず、さらに十分高いエネルギーでは粒子対の生成・消滅が起こりうるため、粒子数を固定したヒルベルト空間 $\mathcal{H}_n$ だけでは状態空間として不十分である。

場の量子論では、まず状態空間としてのヒルベルト空間 $\mathcal{H}$ を与える。自由理論では粒子数が固定されないから、状態空間は
\[
\mathcal{H}=\mathcal{H}_0\oplus \mathcal{H}_1\oplus \mathcal{H}_2\oplus\cdots
\]
のように、$0$ 粒子、$1$ 粒子、$2$ 粒子、$\dots$ の各部分空間をすべて合わせたものになる。ボース粒子では各 $\mathcal{H}_n$ は対称部分空間、フェルミ粒子では反対称部分空間として構成される。後でこの空間をフォック空間と呼び、具体的な式を書き下す。

その上で、時空 $x=(t,\mathbf{x})$ に依存する場を導入する。記法の上では $\hat{\phi}(x), \hat{\psi}(x)$ のように書くが、これは各時空点に数を割り当てる通常の関数ではなく、ヒルベルト空間 $\mathcal{H}$ 上に作用する演算子を時空点で並べたものである。より正確には、滑らかで、ある有限な領域の外では $0$ になる関数 $f$ に対して
\[
\hat{\phi}(f) := \int d^4x\, f(x)\hat{\phi}(x)
\]
のような演算子を与える対象として理解する。以下では記法を簡単にするため、こうした対象を形式的に $\hat{\phi}(x),\hat{\psi}(x)$ と書く。

したがって、場の量子論を定式化する際の基本的なデータは、
\[
\text{状態空間 } \mathcal{H}
\qquad \text{と} \qquad
\text{その上に作用する演算子値分布 } \hat{\phi}(x), \hat{\psi}(x), \dots
\]
である。粒子状態は、これらの場から構成される生成演算子を真空に作用させることで得られるベクトルとして記述される。

以下では、まず最も簡単な例である自由実スカラー場を用いて、この構造を具体的に書き下す。ここで「自由」とは相互作用項がないこと、「スカラー」とはローレンツ変換の下で内部成分を持たないことを意味する。

\subsection{自由スカラー場の古典論}
最も単純な例として、実スカラー場 $\phi(x)$ を考える。以下、この節では見通しをよくするために自然単位系 $c=\hbar=1$ を用いる。この単位系では、特殊相対論の式や量子論の式に現れる $c,\hbar$ をいちいち書かずに済む。ここでも流れは同じであり、まず古典場の運動方程式とエネルギーを定め、その後にそれを量子化する。

場の理論でも、粒子力学と同様に、まず古典的な運動方程式を定めてから量子化する。その出発点になるのがラグランジアン密度 $\mathcal{L}$ である。これは各時空点に割り当てられた量であり、粒子力学のラグランジアン $L$ の場版と思えばよい。これを時空全体で積分した
\[
S = \int d^4x\, \mathcal{L}
\]
を作用という。ここで $d^4x = dt\,d^3x$ は時刻と空間全体にわたる積分を表す。作用が停留するという条件から運動方程式が決まる。

自由スカラー場のラグランジアン密度を
\[
\mathcal{L} = \frac{1}{2}\partial_\mu \phi \, \partial^\mu \phi - \frac{1}{2}m^2\phi^2
\]
とおく。ここで $\partial_\mu$ は時刻と空間に関する微分をまとめた記号であり、$\partial_\mu \partial^\mu$ は時間微分と空間微分をローレンツ対称な形でまとめた演算子である。作用 $S$ に対してオイラー--ラグランジュ方程式を適用すると、
\[
\partial_\mu \partial^\mu \phi + m^2 \phi = 0
\]
が得られる。これはすでに見たクライン--ゴルドン方程式そのものである。したがって、クライン--ゴルドン方程式は「一粒子波動方程式」としてだけでなく、「古典場の運動方程式」としても読めることが分かる。

次に、後でハミルトン形式で量子化するため、場に対する「共役な運動量」を定義する。これは粒子力学における座標 $q$ に対する運動量 $p$ に対応する量であり、
\[
\pi(\mathbf{x},t) := \frac{\partial \mathcal{L}}{\partial \dot{\phi}(\mathbf{x},t)} = \dot{\phi}(\mathbf{x},t)
\]
で与えられる。ここで $\dot{\phi}$ は $\phi$ の時間微分である。

同様に、粒子力学のハミルトニアンに対応する量としてハミルトニアン密度
\[
\mathcal{H} = \pi \dot{\phi} - \mathcal{L}
= \frac{1}{2}\pi^2 + \frac{1}{2}(\nabla \phi)^2 + \frac{1}{2}m^2\phi^2
\]
を得る。全ハミルトニアンはこれを空間全体で積分した
\[
H = \int d^3x\, \mathcal{H}
\]
である。したがって、自由場は「各空間点に自由度をもつ系」、あるいはフーリエ変換すれば「連続無限個の調和振動子の集まり」とみなせる。

\subsection{作用形式と経路積分の位置づけ}
ここまでで作用 $S$ を導入したのは、まず古典場の運動方程式を明示するためである。本節の主な目的は粒子状態と生成・消滅演算子を構成することなので、以下では正準量子化を用いる。一方で、同じ作用から量子論を組み立てる別の方法として、経路積分の考え方もここで位置づけだけを述べておく。

まず粒子力学では、時刻 $t_i$ に位置 $x_i$ から出発し、時刻 $t_f$ に $x_f$ に到達する振幅は
\[
K(x_f,t_f;x_i,t_i)
=
\int_{x(t_i)=x_i}^{x(t_f)=x_f} \mathcal{D}x \, e^{iS[x]}
\]
と書ける。ここで
\[
S[x]=\int_{t_i}^{t_f} dt\, L(x,\dot{x})
\]
は軌道 $x(t)$ に沿った作用であり、$\mathcal{D}x$ は端点を固定したすべての軌道にわたる積分を表す。つまり量子力学では、粒子はただ一つの古典軌道をたどるのではなく、許されるすべての軌道が位相 $e^{iS[x]}$ を伴って干渉し合う。

場の理論では、積分すべき対象が粒子の軌道 $x(t)$ から、時空全体での場の配置 $\phi(x)$ に置き換わる。自由スカラー場なら
\[
Z = \int \mathcal{D}\phi\, e^{iS[\phi]}
\]
を基本量とし、ここで
\[
S[\phi]=\int d^4x \left(
\frac{1}{2}\partial_\mu \phi \, \partial^\mu \phi - \frac{1}{2}m^2\phi^2
\right)
\]
である。積分変数が関数そのものであるため、これはより正確には汎関数積分と呼ばれる。

さらに、場の相関関数は
\[
\langle 0|T\hat{\phi}(x_1)\hat{\phi}(x_2)|0\rangle
=
\frac{1}{Z}\int \mathcal{D}\phi\, \phi(x_1)\phi(x_2)e^{iS[\phi]}
\]
のように表される。左辺の $T$ は、時刻の大きい演算子を左に並べる時間順序積である。したがって、ラグランジアンが与えられれば、そこから相関関数や散乱振幅の計算を系統的に組み立てることができる。

経路積分は、位相 $e^{iS}$ の急速な振動を考えると、作用が停留する配置、すなわち古典方程式を満たす配置の近くが主に寄与する。したがって、この定式化は古典論と量子論のつながりを特に見やすくする。

以下ではまず、粒子状態や生成・消滅演算子の構造が直接見える正準量子化を用いて自由場を記述する。経路積分は、その同じ理論をラグランジアンから組み立てるもう一つの見方として理解しておけばよい。

\subsection{正準量子化}
ここから古典場を量子化する。正準量子化とは、粒子力学で $q,p$ を演算子 $\hat{q},\hat{p}$ に置き換えて交換関係 $[\hat{q},\hat{p}] = i$ を課したのと同じ考え方を、場 $\phi,\pi$ に適用する手続きである。

すなわち、古典場 $\phi,\pi$ を演算子値の場 $\hat{\phi},\hat{\pi}$ に持ち上げ、同時刻での正準交換関係
\[
[\hat{\phi}(\mathbf{x},t), \hat{\pi}(\mathbf{y},t)] = i\delta^{(3)}(\mathbf{x}-\mathbf{y}),
\]
\[
[\hat{\phi}(\mathbf{x},t), \hat{\phi}(\mathbf{y},t)] = 0, \quad
[\hat{\pi}(\mathbf{x},t), \hat{\pi}(\mathbf{y},t)] = 0
\]
を課す。ここで $\delta^{(3)}(\mathbf{x}-\mathbf{y})$ は3次元デルタ関数であり、「異なる空間点の自由度は独立で、同じ点でだけ基本交換関係が成り立つ」ことを表している。

これは、粒子力学における $[\hat{q},\hat{p}] = i$ の連続自由度版である。

\subsection{モード展開と生成・消滅演算子}
自由スカラー場の理論を考える時は、空間の各点ごとに独立な自由度がばらばらにあると考えるよりも、場全体がひとつの力学系を構成し、その振動モードが量子化されると理解すると見通しがよい。
自由場の方程式は線形であるから、その解は平面波の重ね合わせで表せる。量子化した後も同じ考え方が使え、演算子場は各運動量モードの和として
\[
\hat{\phi}(x) =
\int \frac{d^3p}{(2\pi)^3}\frac{1}{\sqrt{2E_{\mathbf{p}}}}
\left(
\hat{a}_{\mathbf{p}} e^{-ip\cdot x}
+ \hat{a}_{\mathbf{p}}^\dagger e^{ip\cdot x}
\right),
\qquad
E_{\mathbf{p}} = \sqrt{\mathbf{p}^2 + m^2}.
\]
と展開できる。これをモード展開という。ここで
\[
p\cdot x := E_{\mathbf{p}}t-\mathbf{p}\cdot\mathbf{x}
\]
である。指数関数 $e^{\mp ip\cdot x}$ は各運動量モードの時空依存性を表し、$E_{\mathbf{p}} = \sqrt{\mathbf{p}^2 + m^2}$ は相対論的エネルギー・運動量関係である。

ここで $\hat{a}_{\mathbf{p}}$ と $\hat{a}_{\mathbf{p}}^\dagger$ は、それぞれ運動量 $\mathbf{p}$ をもつ量子を一つ消滅・生成する演算子であり、
\[
[\hat{a}_{\mathbf{p}}, \hat{a}_{\mathbf{q}}^\dagger] = (2\pi)^3 \delta^{(3)}(\mathbf{p}-\mathbf{q}),
\quad
[\hat{a}_{\mathbf{p}}, \hat{a}_{\mathbf{q}}] = 0
\]
を満たす。この交換関係は、各運動量モードが調和振動子と同じ代数構造を持つことを示している。

ハミルトニアンはこれらを用いて書き直せる。すると各モードは「エネルギー $E_{\mathbf{p}}$ をもつ量子が何個あるか」で記述される形になる。真空にも形式的には無限大の零点エネルギーが現れるが、まずはその基準値を引いた量を見るため、生成演算子を左、消滅演算子を右に並べる正規順序をとると、
\[
:\hat{H}:=
\int \frac{d^3p}{(2\pi)^3} E_{\mathbf{p}} \,
\hat{a}_{\mathbf{p}}^\dagger \hat{a}_{\mathbf{p}}
\]
となる。したがって、一つ一つの粒子は「場の一つのモードを一量子だけ励起した状態」に対応している。

\subsection{フォック空間と多粒子状態}
以上のモード展開と交換関係から、自由実スカラー場の一粒子空間は、いま採用している規格化のもとで
\[
\mathfrak{h}_\phi
:=
L^2\!\left(\mathbb{R}^3,\frac{d^3p}{(2\pi)^3}\right)
\]
ととる。この選択は、運動量波束 $f(\mathbf{p})$ に対して
\[
|f\rangle
:=
\int \frac{d^3p}{(2\pi)^3} f(\mathbf{p}) \hat{a}_{\mathbf{p}}^\dagger |0\rangle
\]
とおけば、
\[
\langle f|g\rangle
=
\int \frac{d^3p}{(2\pi)^3} f(\mathbf{p})^* g(\mathbf{p})
\]
が成り立つ。

したがって、自由実スカラー場のヒルベルト空間は、この一粒子空間の上に作ったボースフォック空間
\[
\mathcal{H}_\phi
=
\mathcal{F}_+(\mathfrak{h}_\phi)
=
\bigoplus_{n=0}^{\infty}\operatorname{Sym}^n \mathfrak{h}_\phi
\]
である。真空状態 $|0\rangle$ は $n=0$ 粒子部分空間に属するベクトルであり、生成演算子を作用させることで高粒子数部分空間へ移る。

粒子がひとつも存在しない状態を真空状態と呼び、$|0\rangle$ と書く。これは
\[
\hat{a}_{\mathbf{p}}|0\rangle = 0 \qquad (\forall \mathbf{p})
\]
で定義される。つまり、どの運動量モードについても「これ以上消せる粒子がない」状態である。

このとき、一粒子状態は $\hat{a}_{\mathbf{p}}^\dagger|0\rangle$、二粒子状態は $\hat{a}_{\mathbf{p}}^\dagger \hat{a}_{\mathbf{q}}^\dagger|0\rangle$ のように作られる。一般に
\[
|\mathbf{p}_1,\dots,\mathbf{p}_n\rangle
= \hat{a}_{\mathbf{p}_1}^\dagger \cdots \hat{a}_{\mathbf{p}_n}^\dagger |0\rangle
\]
が $n$ 粒子状態を与える。ここでは、以前に固定した粒子数 $n$ を最初から決めておく必要はなく、生成演算子を何回作用させたかによって粒子数が決まる。

このように、場の量子論で用いるヒルベルト空間は、粒子数が固定されたひとつの空間ではなく、$0$ 粒子、$1$ 粒子、$2$ 粒子、$\dots$ の各部分空間をすべて合わせた空間である。これをフォック空間という。相互作用を導入すると、時間発展の中でこれら異なる粒子数の部分空間が混ざり合う。したがって、粒子対生成や消滅は、異なる粒子数部分空間を結ぶ時間発展として記述される。

\subsection{ディラック場への一般化}
同じ考え方はスカラー場だけでなく、電子のようなスピン $1/2$ 粒子を記述するディラック場にも適用できる。ディラック方程式も、もはや一粒子の波動方程式としてではなく、スピノール場 $\psi(x)$ に対する古典場の方程式として読み直すのが本質的である。

ディラック場では、一粒子空間として電子と陽電子を別々にとる。運動量表示では
\[
\mathfrak{h}_e
=
\mathfrak{h}_{\bar e}
:=
L^2\!\left(\mathbb{R}^3,\frac{d^3p}{(2\pi)^3};\mathbb{C}^2\right)
\]
とする。ここで $\mathbb{C}^2$ は各運動量に対してスピン自由度が2つあることを表している。

したがって、ディラック場のヒルベルト空間はフェルミフォック空間
\[
\mathcal{H}_{\psi}
=
\mathcal{F}_-(\mathfrak{h}_e\oplus \mathfrak{h}_{\bar e})
\]
である。直和 $\mathfrak{h}_e\oplus \mathfrak{h}_{\bar e}$ は、電子と陽電子の一粒子自由度をまとめている。

ただし、電子はフェルミ粒子であるから、ボース粒子のような交換関係ではなく、同時刻の反交換関係
\[
\{\hat{\psi}_\alpha(\mathbf{x},t), \hat{\psi}_\beta^\dagger(\mathbf{y},t)\}
= \delta_{\alpha\beta}\delta^{(3)}(\mathbf{x}-\mathbf{y})
\]
を課す。ここで添字 $\alpha,\beta$ はスピノールの成分を表す。これにより、電子のようなフェルミ粒子に対して同じ量子状態を二つ以上占有できないというパウリの排他原理が自動的に組み込まれる。

ディラック場でもモード展開を行うと、電子を生成する演算子と陽電子を生成する演算子が別々に現れる。したがって、反粒子の存在はもはや後付けの再解釈ではなく、場の量子化そのものの帰結となる。これは先に見た「負エネルギー解をどう解釈するか」という問題に対する、より根本的な答えである。

\subsection{相互作用への道筋}
自由場の量子化によって、状態空間と生成・消滅演算子の構造は得られた。次の課題は、異なる場をどのような原理で結合するかである。その最も基本的な例が、電磁場 $A_\mu(x)$ とディラック場 $\psi(x)$ の相互作用を記述する量子電磁力学（QED）である。以下では、自由理論から出発し、対称性の要請だけで相互作用項がどのように定まるかを見る。
